\SourceLeaf{335}%
\Article{34}
Von den sieben Koeffizienten (1), (2) usw. ist der letzte konstant. Es gilt
\[
\log(7)=7{,}6287228032(-20)
\]
und
\[
\log3(7)=8{,}1058440580(-20)
\]
Sobald man feste Werte für die Abmessungen des Ellipsoids gewählt hat, werden die übrigen Koeffizienten zu Funktionen der Breite \(B\). Man kann sie also in einer Tafel mit \(B\) als Argument zusammenstellen. Steht eine solche Tafel zur Verfügung, so ist die Rechnung nach dieser Methode für das Ellipsoid ebenso bequem wie für die Kugel.

Am Schluss dieser Abhandlung füge ich eine solche Tafel für die Zone von \(51^0\) bis \(54^0\) bei, in der die Werte von \(B\) in Schritten von einer Minute fortschreiten. Dazu bemerke ich Folgendes.

Die Tafel hängt von den Parametern des Ellipsoids nur insofern ab, als ihr eine bestimmte Abplattung oder ein bestimmter Wert von \(e\) zugrunde liegt. Es ist der Wert, den die letzte von \textsc{Bessel} ausgeführte Rechnung ergeben hat und der auch der Tafel zur ersten Abhandlung zugrunde liegt (siehe Art. 5). Damit der Zahlenwert von \(a\) nur von der Abplattung abhängt, dient als Einheit anstelle der Toise oder eines anderen willkürlichen Maßes der zehnmillionste Teil des Erdmeridians. Damit lässt sich \(a\) unmittelbar aus \(e\) bestimmen, und zwar mithilfe der Gleichung
\[
\begin{aligned}
\pi a(1-\tfrac14ee-\frac{1\mathord{.}3}{4\mathord{.}16}e^4-\frac{1\mathord{.}3\mathord{.}15}{4\mathord{.}16\mathord{.}36}e^6\\
{}-\frac{1\mathord{.}3\mathord{.}15\mathord{.}35}{4\mathord{.}16\mathord{.}36\mathord{.}64}e^8
-\frac{1\mathord{.}3\mathord{.}15\mathord{.}35\mathord{.}63}{4\mathord{.}16\mathord{.}36\mathord{.}64\mathord{.}100}e^{10}-\text{usw.})\\
=20000000
\end{aligned}
\]
deren Bildungsgesetz offensichtlich ist, oder mithilfe der gleichwertigen Gleichung
\[
a=\frac{20000000}{\pi}(1+\tfrac14ee+\tfrac7{64}e^4+\tfrac{15}{256}e^6+\tfrac{579}{16384}\mathord{.}e^8+\tfrac{1515}{65536}e^{10}+\text{usw.})
\]
Mit jenem Wert von \(e\) erhält man so die folgenden Ergebnisse.
\[
\begin{aligned}
a&=6376851{,}447\\
\log a&=6{,}8046062999
\end{aligned}
\]
Bei der Anwendung unserer Tafel muss selbstverständlich auch \(r\) zunächst in derselben Einheit ausgedrückt werden. Für diese Umrechnung muss man, ausgehend von dem Wert, den
